\section{Solving the two-cell response, not merely one Euler step}
The earlier example froze the initial flux before taking a finite step.
We can also solve its continuous response exactly, which helps explain
why those are different objects. Keep total mass twenty, references
$(10,10)$, scales two, one forward edge, mobility $1/2$, zero threshold
and no external drive. Write
\begin{equation}
 x_1=10+y,\qquad x_2=10-y,\qquad y(0)=8.
\end{equation}
For $y\geq0$, the force is $y/2$ and the flux is $y/4$.
Conservation gives $\dot y=-y/4$. Separate variables, or differentiate
the proposed solution, to obtain
\begin{equation}
 y(t)=8e^{-t/4},\qquad V(t)=y(t)^2/4=16e^{-t/2}.
\end{equation}
For every finite nonnegative time the stocks remain nonnegative and the
force stays in the chosen forward direction. The reference is approached
asymptotically. It is not reached at the finite duration four simply
because the frozen Euler step of that duration landed there.

The solution is a conditional analytic example, not an experimental run.
It depends on the chosen continuous response law. Replacing that law by
random finite actions does not preserve the differential equation, so the
exponential expression cannot be advertised as the predicted recovery
curve of the new programme.

\section{A response threshold can leave residual deviation}
Keep the same example but take a fixed threshold $0<\theta<4$.
While $y>2\theta$, the forward flux is
$(y-2\theta)/4$. Therefore
\begin{equation}
 y(t)=2\theta+(8-2\theta)e^{-t/4}.
\end{equation}
This tends to $2\theta$, not zero. The limiting potential is
$\theta^2$. Below the activation threshold the directed law has no
forward flow. If the initial $y$ were already at most $2\theta$ and
nonnegative, this edge alone would remain inactive.

This elementary calculation shows why nonincrease of a potential is not
the same as attraction to its minimum. The threshold is a response
parameter, not an Allee stock threshold. Removing the latter from the
Gaussian geometry does not remove the former from an optional controller
that explicitly declares it. The random-affordable proposal contains
neither response threshold by default.

\section{What positive clipping preserves and what it changes}
Suppose an otherwise admissible historical flux is reduced edge by edge
by factors $0\leq\alpha_e\leq1$, so the accepted rate is
$\widetilde J_e=\alpha_eJ_e$. On every active edge the original law
has $f_e\geq\theta_e\geq0$. Thus
$f_e\widetilde J_e\geq0$, and with no drive the chain rule still gives
\begin{equation}
 \dot V=-\sum_e f_e\widetilde J_e\leq0.
\end{equation}
This sign argument is weaker than the original equality expressed in
$J_e^2/M_e+\theta_eJ_e$. The accepted rate no longer necessarily
satisfies the original constitutive equation with the old mobility.
Preserving a sign does not preserve every quantitative statement.

Nor does this argument cover an arbitrary resolver. A resolver that changes
directions, introduces compensating transformations or changes the state
before evaluation needs its own analysis. In discrete time even positive
clipping still leaves a curvature term to check for the accepted increment.
A feasibility improvement is not automatically a numerical descent proof.

\section{The three meanings of a safe step}
A step can be safe with respect to nonnegative stock, safe with respect
to a specified reserve certificate, or safe with respect to nonincrease
of the chosen potential. These properties can overlap, but they are not
synonyms. The two-cell duration-nine example earlier was physically
nonnegative and nevertheless increased potential. A tiny step with the
wrong physical support can be impossible however small its numerical
truncation error would be.

For this reason the detailed historical proofs that follow should be
read as several conditional certificates, not one universal safety
theorem. Their line-by-line derivations remain valuable precisely because
the reader can see where each premise enters and decide whether it belongs
to the current model.
